%\chapter{Vibrations of Membranes}\label{vibrations-of-membranes}
193. The theoretical membrane is a perfectly flexible and infinitely thin
lamina of solid matter, of uniform material and thickness, which is stretched
in all directions by a tension so great as to remain sensibly unaltered
during the vibrations and displacements contemplated. If an imaginary line be
drawn across the membrane in any direction, the mutual action between the two
portions separated by an element of the line is proportional to the length of
the element and perpendicular to its direction. If the force in question be
\(T_1 ds\), \(T_1\) may be called the \textit{tension of the membrane}, it is
a quantity of one dimension in mass and \(-2\) in time.

The principal problem in connection with this subject is the investigation of
the transverse vibrations of membranes of different shapes, whose boundaries
are fixed. Other questions indeed may be proposed, but they are of
comparatively little interest ; and, moreover, the methods proper for solving
them will be sufficiently illustrated in other parts of this work. We may
therefore proceed at once to the consideration of a membrane stretched over
the area included within a fixed, closed, plane boundary.

% LaTeX page 175
\setcounter{equation}{0}
194. Taking the plane of the boundary as that of \(x y\), let \(w\)
denote the small displacement therefrom of any point \(P\)
of the membrane. Round \(P\) take a small area \(S\),
and consider the forces acting upon it parallel to \(z\).
The resolved part of the tension is expressed by
\[ T_1 \int \frac{\del w}{\del n} ds,\]
where \(ds\) denotes an element of the boundary of \(S\), and \(d \vec{n}\) an
element of the normal to the curve drawn outwards. This is balanced by
the reaction against acceleration measured by \(\rho S \ddot{y} \),
\(\rho\), being a symbol of one dimension in mass and \(-2\) in length
denoting the superficial density. Now by Green's theorem, if
\[
\Grad = \left(\frac{\del}{\del x}, \frac{\del}{\del y} \right), \quad
\Lap = \frac{\del^2}{\del x^2} + \frac{\del^2}{\del y^2},
\]
\[\int \frac{\del w}{\del n}ds=\iint_D \Lap w \, dS \approx \Lap w \,|S|\]
ultimately, and thus the equation of motion is (1)
\begin{equation} \ddot{w} = \frac{T_1}{\rho} \Lap w.\end{equation}

The condition to be satisfied at the boundary is of course \(w=0\).

The differential equation may also be investigated from the expression for the potential energy, which is found by multiplying the tension by the superficial stretching.  The altered area is
\[ \iint \sqrt{ 1 + {\norm{\Grad w}}^2} \, dx \, dy ;\]
from which V is easily found by an integration by parts, and thus
\begin{equation}
V=\frac{T_1}{2} \iint {\norm{\Grad w}}^2 \, dx \, dy;
\end{equation}
from which \(\delta V\) is easily found by an integration by parts.

If we write
\(T_1/\rho = c^2\),
then \(c\) is of the nature of a velocity, and the
differential equation is
\begin{equation}\label{equation194_3} \ddot{w} = c^2 \Lap w.\end{equation}

\setcounter{equation}{0}
\section{Levy solution part 1, 195}

195. We shall now suppose that the boundary of the membrane is the rectangle
formed by the coordinate axes and the lines \(x=a\),\(y=b\). For every point
within the area, equation~(\ref{equation194_3}) is satisfied, and for every
point on the boundary \(w=0\).

A particular integral is evidently
\begin{align}
w&=\sin \left(\frac{m \pi x}{a}\right) \sin \left(\frac{n \pi y}{b}\right) \cos(pt),
p^2 &= c^2 \pi^2 \left(\frac{m^2}{a^2} + \frac{n^2}{b^2} \right),
\end{align}
and \(m\) and \(n\) are integers; and from this the general solution may be derived. Thus (3)
\begin{equation}
w=\sum_{m,n \geq 1}
\sin \left(\frac{m \pi x}{a}\right) \sin \left(\frac{n \pi y}{b}\right)
(A_{mn} \cos(pt) + B_{mn} \sin(pt)).
\end{equation}

That this result is really general may be proved \textit{a posteriori}, by
shewing that it may be adapted to express arbitrary initial circumstances.

Whatever function of the coordinates \(w\) may be, it
can be expressed for \(0 \leq x \leq a\) by the series
\[\sum_m Y_m \sin \left(\frac{m \pi x}{a}\right),\]
where the coefficients \({\{ Y_m \}}_{m \geq 1}\) are independent of \(x\).
Again whatever function of \(y\) and any one of the coefficients \(Y_m\) be,
it can be expanded, \(0 \leq y \leq b\) in the series
\[ \sum_n C_n \sin \left(\frac{n \pi y}{b}\right),\]
\({\{ C_n \}}_{n \geq 1}\) are constants.
From this we conclude that any function of
\(x\) and \(y\) can be expressed within the limits of the rectangle by the
double series
\[
\sum_{m,n \geq 1} A_{mn}
\sin \left(\frac{m \pi x}{a}\right) \sin \left(\frac{n \pi y}{b}\right)
\]
and therefore that the expression for \(w\) in (3) can be adapted to
arbitrary initial values of \(w\) and \(\dot{w}\).
For the rectangle \(D = [0,a] \times [0,b]\),
\begin{equation}
(A_{mn}, B_{mn}) = \frac{4}{ab} \iint_D \left(w_o, \frac{ \dot{w}_o }{p} \right)
\sin \left(\frac{m \pi x}{a}\right) \sin \left(\frac{n \pi y}{b}\right) \, dx \, dy
\end{equation}

The character of the normal functions of a given rectangle,
\[
\sin \left(\frac{m \pi x}{a}\right) \sin \left(\frac{n \pi y}{b}\right)
\]
as depending on \(m\) and \(n\), is easily understood. If
\(m\) and \(n\) be both unity, \(w\)
retains the same sign over the whole of the rectangle, vanishing at
the edge only; but in any other case there are nodal lines running
parallel to the axes of coordinates. The number of the nodal lines
parallel to \(x\) is \(n-1\), their equations being
\[ y=k b/n,\quad 1 \leq k < n.\]

In the same way the equations of the nodal lines parallel to
\(y\) are
\[ x = k a/m, \quad 1 \leq k < m,\]
being \(m-1\), in number. The nodal system divides the
rectangle into mn equal parts, in each of which the numerical value of
\(w\) is repeated.

%\[ w = \sum_{m,n \geq 1}
% \sin \left(\frac{m \pi x}{a}\right) \sin \left(\frac{n \pi y}{b}\right)
% V=\frac{T_1}{2} \iint {\norm{\Grad w}}^2 \, dx \, dy;
%\iint_D f(x,y) \, dx \, dy
% xxxx

\setcounter{equation}{0}
\section{Levy solution part 2, 196}
196. The expression for w in terms of the normal functions is

\[w= Xz nq sin \sin(Y) (1),\]
where c, etc. are the normal coordinates. We proceed to form the
expression for V in terms of c,,: We have

\[we) = 0 254.7 00s Ms \sin()\]

\[\sin()= \cos(ay)\]

In integrating these expressions over the area of the rectangle the
products of the normal coordinates disappear, and we find

(2),
the summation being extended to all integral values of m and n.

The expression for the kinetic energy is proved in the same way to be
pab
team (3), from which we deduce as the normal equation of motion

(4).

In this equation

apo. mre. mry ,= "f Zin (5 dx dy (5),
if Z dx dy denote the transverse force acting on the element dx dy.

Let us suppose that the initial condition is one of rest under the
operation of a constant force Z, such as may be supposed to arise from
gaseous pressure. At the time c=0, the impressed force is removed, and
the membrane left to itself. Initially the equation of equilibrium is

\(c () = ab \Phi\) (6)
whence () o is to be found. The position of the system at time t is then
given by

(7),
in conjunction with (1).

In order to express \(\Phi\) , we have merely to substitute for Z its value
in (5), or in this case simply to remove Z from under the integral sign.
Thus

\(\Phi = Z () () dx dy\)

= Z a b (1 - (mr) (1 - (nz).

We conclude that \(\Phi\) vanishes, unless m and n are both odd, and that
then

\(\Phi = 4ab/() Z.\)

Accordingly, m and n being both odd,= (8),

where p 2 = (9).

This is an example of (8), Article 101.

If the membrane, previously at rest in its position of equilibrium, be
set in motion by a blow applied at the point , the solution is

= 4/abp (2) (8) fd dy (pt)(10).

[As an example of forced vibrations, suppose that a harmonic
force acts at the centre. Unless \(m\) and \(n\) are both odd,
\( \Phi_{mn} = 0\), and in the case reserved
\[ \Phi_{mn} = \pm Z_1 \cos (qt) \]
(11), and in the case reserved
where
\( Z_1\)
is the whole force acting at time \(t\), and
\(\pm\) represents
\( \sin (m \pi/2) \sin (n \pi/2)\). From (4) and (9) we have
\[ \phi_{mn} = \pm 4 Z_1 \cos (qt).  \]
(12) In the case of a square membrane, \(p\) is a symmetrical function
of \(m\) and \(n\).  When \(m\) and \(n\) are unequal, the terms occur in pairs,
such as
\[ \left( \sin (m \pi x/a) \sin (n \pi y/a) + \sin (n \pi x/a) \sin (m \pi y/a) \right),\]
a combination symmetrical as between x and y. The vibration is
of course similarly related as well to the four sides as to the four
comers of the square.

In the neighbourhood of the centre, where the force is applied,
the series loses its convergency, and the displacement w tends to
become (logarithmically) infinite.]

\setcounter{equation}{0}
\section{Degenerate modes of square membranes, 197}

197. The frequency of the natural vibrations is found by ascribing
different integral values to m and n in the expression

(1).

For a given mode of vibration the pitch falls when either side of the
rectangle is increased. In the case of the gravest mode, when m=1, n=1,
additions to the shorter side are the more effective; and when the form
is very elongated, additions to the longer side are almost without
effect.

When a2 and b2 are incommensurable, no two pairs of values of m and n
can give the same frequency, and each fundamental mode of vibration has
its own characteristic period. But when a2 and b2 are commensurable, two
or more fundamental modes may have the same periodic time, and may then
coexist in any proportions, while the motion still retains its simple
harmonic character. In such cases the specification of the period does
not completely determine the type. The full consideration of the problem
now presenting itself requires the aid of the theory of numbers; but it
will be sufficient for the purposes of this work to consider a few of
the simpler cases, which arise when the membrane is square. The reader
will find fuller information in Riemann's lectures on partial
differential equations.

If a=),

The lowest tone is found by putting m and n equal to unity, which gives
only one fundamental mode :

W = (BID) (Pb os (3).

Next suppose that one of the numbers m, n is equal to 2, and the other
to unity. In this way two distinct types of vibration are obtained,
whose periods are the same. If the two vibrations be synchronous in
phase, the whole motion is expressed by

w= C (2) (724) D (2) (2H) (pt. (4);
so that, although every part vibrates synchronously with a harmonic
motion, the type of vibration is to some extent arbitrary,

Four particular cases may be especially noted. First, if D=0, (5),
which indicates a vibration with one node along the line x= 4a.
Similarly if C=0, we have a node parallel to the other pair of edges.
Next, however, suppose that C and D are finite and equal, then w is
proportional to

\[\sin(27)  x  + \sin(2) \sin(2),\]
which may be put into the form

\[2 \sin(7) \sin(7) \cos(8) ,\]

This expression vanishes, when \(\sin(7) =0,\) or \(\sin(24) =0\) or
again, when

\[\cos()+ \cos(4) =0,\]

The first two equations give the edges, which were originally
assumed to be nodal; while the third gives y+ a= a, representing one
diagonal of the square. - In the fourth case, when C=-D, we obtain for
the nodal lines, the edges of the square together with the diagonal y=.
The figures represent the four cases.

Fig. 32, D=0. C+D=0.
ANY

[Frequency (referred to gravest) = 1‘58.]

For other relative values of C and D the interior nodal line is curved,
but is always analytically expressed by

and may be easily constructed with the help of a table of logarithmic
cosines.

[Frequency = 2 00.]


The next case in order of pitch occurs when m=2, n=2, The values of m
and n being equal, no alteration is caused by their interchange, while
no other pair of values gives the same frequency of vibration. The only
type to be considered is accordingly

\[w = \sin(m) \sin(m) \cos(pt),\]
whose nodes, determined by the equation

\[n = \sin(7) \cos(t)\]
are (in addition to the edges) the straight lines a= ta y=te. Fig. 33,

The next case which we shall consider is obtained by ascribing to m, n
the values 3, 1, and 1, 3 successively. We have

w= C (a) ()+ D () (zu) (pt). The nodes are given by () (7Y) 0 (4c08' -
1) + D 4 008 - -i=0, a or, if we reject the first two factors, which
correspond to the edges, (4 (t) + D (4 (Y 1)= 0 (7).

If C=0,we have y=a, y=

If D=0, x=fa, o=.

If C=-D () = + (74, a a whence, y=X,

which represent the two diagonals. Lastly, if C= D, the equation of the
node is (e) + (a) =b or 1+ (7) + (2) = H (B),

Fig. 84. C+D=0.

= () on

[Frequency = 2 24.]

In case (4) when z2=a, y=4a, or Za; and similarly when =a, w= a, or Z.
Thus one half of each of the lines joining the middle points of opposite
edges is intercepted by the curve.

[The diameters of the nodal curve parallel to the sides of the
square are thus equal to fa- Those measured along the diagonals
are sensibly smaller, equal to f . a, or 471 a.]


It should be noticed that in whatever ratio to one another C and D may
be taken, the nodal curve always passes through the four points of
intersection of the nodal lines of the first two cases, C=0, D=0. If the
vibrations of these cases be compounded with corresponding phases, it is
evident that in the shaded compartments of Fig. (35) the directions of
displacement

Fig. 35. Ki


are the same, and that therefore no part of the nodal curve is to be
found there; whatever the ratio of amplitudes, the curve must be drawn
through the unshaded portions, When on the other hand the phases are
opposed, the nodal curve will pass exclusively through the shaded
portions.

When m=3, n=3, the nodes are the straight lines parallel to the edges
shewn in Fig. (36).

Fig. 36.


The last case
m=2, n=3. The nodal system is

[Frequency - 3.OO.]

C (2) gin 27Y 4 D (27) gin 27Y - 0 a a a a
or, if the factors corresponding to the edges be rejected, o(4 (t) (7Y)
+ D = (4 (t)) (9).

If C or vanish, we fall back on the nodal systems of the component
vibrations, consisting of straight lines parallel to the edges. If C=D,
our equation may be written

( (7 +-cos wy) (4 008 = a 008 TH 1) = 0. (10),

of which the first factor represents the diagonal y+a=a, and the second
a hyperbolic curve.

If C=- D, we obtain the same figure relatively to the other diagonal'.

\setcounter{equation}{0}
\section{Maybe more Degenerate modes of qquare membranes, 198}

198. The pitch of the natural modes of a square membrane, which is
nearly, but not quite uniform, may be investigated by the genera method
of Article 90.

We will suppose in the first place that m and n are equal. In this case,
when the pitch of a uniform membrane is given, the mode of its vibration
is completely determined. If we now conceive a variation of density to
ensue, the natural' type of vibration is in general modified, but the
period may: be calculated approximately without allowance for the change
of type.

Dangling\footnote{ Lame, Crelle, sur V élasticité, p. 129.}

We have T= 4 o. + p) Pant (a) () = dx dy

yq sf 9 Mme . 2 mTy = Down +f () () dx dy of which the second term is
the increment of 7' due to 8p. Hence
if w cc (pt), and P denote the value of p previously to variation,
we have Pot Pmt=1-4 al. = (MO ( dr dy(V),
where Pam = 2mm , and c=T7,+p, For example, if there be a small load M
attached to the middle of the square,
(2),
in which (4m) vanishes, if m be even, and is equal to unity, if m be
odd. In the former case the centre is on the nodal line of the unloaded
membrane, and thus the addition of the load produces no result.

When, however, m and n are unequal, the problem, though remaining
subject to the same general principles, presents a peculiarity different
from anything we have hitherto met with. The natural type for the
unloaded membrane corresponding to a specified period is now to some
extent arbitrary; but the introduction of the load will in general
remove the indeterminate element. In attempting to calculate the period
on the assumption of the undisturbed type, the question will arise how
the selection of the undisturbed type is to be made, seeing that there
are an indefinite number, which in the uniform condition of the membrane
give identical periods. The answer is that those types must be chosen
which differ infinitely little from the actual types assumed under the
operation of the load, and such a type will be known by the criterion of
its making the period calculated from it a maximum or minimum.

As a simple example, let us suppose that a small load M-is attached to
the membrane at a point lying on the line 2 = 4a, and that we wish to
know what periods are to be substituted for the two equal periods of the
unloaded membrane, found by making

m=1,n=2, or m=2,n=1. It is clear that the normal types to be chosen, are
those whose nodes are represented in the first two cases of Fig. (82).
In the first case the increase in the period due to the load is zero,
which is the least that it can be; and in the second case the increase
is the greatest possible. If 8 be the ordinate of M, the kinetic energy
is altered in the ratio

2 2 fs : boy O (t)228, and thus pat: Pte 1 eM int BOE (3)


while Pu' = P,; =P a The ratio characteristic of the interval between
the two natural ' tones of the loaded membrane is thus approximately

1 a n t TE ca (4), ap a

If 8 = fa, neither period is affected by the load.

As another example, the case, where the values of m and n are 3 and 1,
considered in Article 197, may be referred to. With a load in the
middle, the two normal types to be selected are those corresponding to
the last two cases of Fig. (84), in the former of which the load has no
effect on the period.

The problem of determining the vibration of a square membrane which
carries a relatively heavy load is more difficult, and we shall not
attempt its solution. But it may be worth while to recall to memory the
fact that the actual period is greater than any that can be calculated
from a hypothetical type, which differs from the actual one.

199. The preceding theory of square membranes includes a good deal more
than was at first intended. Whenever in a vibrating system certain parts
remain at rest, they may be supposed to be absolutely fixed, and we thus
obtain solutions of other questions than those originally proposed. For
example, in the present case, wherever a diagonal of the square is
nodal, we obtain a solution applicable to a membrane whose fixed
boundary is an isosceles right-angled triangle. Moreover, any mode of
vibration possible to the triangle corresponds to some natural mode of
the square, as may be seen by supposing two triangles put together, the
vibrations being equal and opposite at points which are the images of
each other in the common hypotenuse. Under these circumstances it is
evident that the hypotenuse would remain at rest without constraint, and
therefore the vibration in question is included among those of which a
complete square is capable.

The frequency of the gravest tone of the triangle is found by putting m
=1, n=2 in the formula

(1),

0/5

and is therefore equal to Sa

The next tone occurs, when m=3,n=1. In this case

p ct0 Foo (2), Fig. 37

as might also be seen by noticing that the triangle divides itself into
two, Fig. (87), whose sides are less than those of the whole triangle in
the ratio 72 : 1.

For the theory of the vibrations of a membrane whose boundary is in the
form of an equilateral triangle, the reader is referred to Lamé's
`Lecons sur l'élasticité'' It is proved that the frequency of the
gravest tone is c+h, where h is the height of the triangle, which is the
same as the frequency of the gravest tone of a square whose diagonal is h.

% Many articles refer back to this one.
200. When the fixed boundary of the membrane is circular, the first step
towards a solution of the problem is the expression of the general
differential equation in polar coordinates. This may be effected
analytically; but it is simpler to form the polar equation de novo by
considering the forces which act on the polar element of area rd dr. As
in Article 194 the force of restitution acting on a small area of the
membrane is . dw d (dw d (dw

- 7, 92 de=- 1,15, (Ge 748) ar-+ 5 (Sy ar) as Pw idw 1dw =- 1, .(Te + et
a apt and thus, if 7, + p =c' as before, the equation of motion is fw
(dw idw idw ene sts ae (1).

The subsidiary condition to be satisfied at the boundary is that w=0,
when r=a.

In order to investigate the normal component vibrations we have now to
assume that w is a harmonic function of the time. Thus, if wo (pt-e),
and for the sake of brevity we write p+c=k, the differential equation
appears in the form (2), in which is the reciprocal of a linear
quantity, Now whatever may be the nature of w as a function of r and 0,
it can be expanded in Fourier's series w = w, + w (8+a) + w (2) (8 +4,)
+(8),

in which w,, w,, etc. are functions of r, but not of 8. The result of
substituting from (3) in (2) may be written iw, , 1 dw, n 2 oe ode + e
Ty) wa (m (0+ a,)=0,

the summation extending to all integral values of n If we multiply this
equation by (4), and integrate with respect to 8 between the limits 0
and 27, we see that each term must vanish separately, and we thus obtain
to determine w, as a function of r

Pu, Ldv, (e-3)w, = (A),

in which it is a matter of indifference whether the factor (n)(0 +4,) be
supposed to be included in w, or not.

The solution of (4) involves two distinct functions of 1, each
multiplied by an arbitrary constant. But one of these functions becomes
infinite when r vanishes, and the corresponding particular solution must
be excluded as not satisfying the prescribed conditions at the origin of
coordinates. This point may be illustrated by a reference to the simpler
equation derived from (4) by making and n vanish, when the solution in
question reduces to w=log r, which, however, does not at the origin
satisfy \(\Grad 2 w =0\), as may be seen from the value of da,
integrated round a small circle with the origin for centre. In like
manner the complete integral of (4) is too general for our present
purpose, since it covers the case in which the centre of the membrane is
subjected to an external force.

[transcribers note.  The discussion of Bessel's functions (see DLMF)
in section 200 only is condensed.]

The other function of 1, which satisfies (4), is the Bessel's function
of the first kind and the nth order, denoted by \(J_n(r)\).
The Frobenius method applied to Bessel's equation is a way to derive the
series expansion of \(J_n(z)\) about \(z=0\).
Next the recursion relations for \(J_n(z)\) and \( d J_n/dz\) can be derived.
Bessel's integrals are the original form.
From the integral form, it is evident that
\(J_n\),
and its differential coefficients with respect to \(z\) are always less than unity.
\todo[inline]{what is a differential coefficient with respect to \(z\) ?}

The series expansion about \(\infty\) is related to the asymptotic expansions
of the Bessel functions.  And the form of these expansions suggests the
spherical Bessel functions, which can be written using Rayleigh's formulas
(see A and S).

201. In accordance with the notation for Bessel's functions the
expression for a normal component vibration may therefore be written
w=PJ,(r) (n) (pt+e) and the boundary condition requires that FT (40) =0.
(2),
an equation whose roots give the admissible values of , and therefore of

The complete expression for w is obtained by combining the particular
solutions embodied in (1) with all admissible values of and n, and is
necessarily general enough to cover any initial circumstances that may
be imagined. We conclude that any function of r and may be expanded
within the limits of the circle r=a in the series

\[w= Z fp \cos(n) + \sin(n)\]

(3). For every integral value of n there are a series of values of ,
given by (2); and for each of these the constants c and are arbitrary.

The determination of the constants is effected in the usual way. Since
the energy of the motion is equal to

bp a (4),

and when expressed by means of the normal coordinates can only involve
their squares, it follows that the product of any two of the terms in
(3) vanishes, when integrated over the area of the circle. Thus, if we
multiply (3) by J, (r) (n), and integrate, we find

=c (J, (er) (nd) r dr dO

= b.0 i F, (e0) (5)

by which c is determined. The corresponding formula for is obtained by
writing (n for (n), A method of evaluating the integral on the right
will be given presently. Since c and each contain two terms, one varying
as (pt) and the other as (pt), it is now evident how the solution may be
adapted so as to agree with arbitrary initial values of w and w.


202. Let us now examine more particularly the character of the
fundamental vibrations. If n=0, w is a function of r only, that is to
say, the motion is symmetrical with respect to the centre of the
membrane. The nodes, if any, are the concentric circles, whose equation
is

J, (r) Oe (1).

When n has an integral value different from zero, w is a function of 6
as well as of r, and the equation of the nodal system takes the form

FT, (eer) (2) (8 - a) =0... (2).

The nodal system is thus divisible into two parts, the first consisting
of the concentric circles represented by

A (29 ed (8), and the second of the diameters Tv B= at (2M +1) Fr (4),

where m is an integer. These diameters are in number, and are ranged
uniformly round the centre; in other respects their position is
arbitrary. The radii of the circular nodes will be investigated further
on.

203. The important integral formula s, (er) Jy (xr) rir = 0 (1),

where x and ' are different roots of S,, (KA) =0 (Q),

may be verified analytically by means of the differential equations
satisfied by J, (r), J,(r); but it is both simpler and more instructive
to begin with the more general problem, where the boundary of the
membrane is not restricted to be circular,

The variational equation of motion is

V + w w dx dy =0 (3)

where

V= (1/2) T1 dx dy (4),

and therefore

V (1/2) T1 dx dy (5),

In these equations w refers to the actual motion, and w to a
hypothetical displacement consistent with the conditions to which the
system is subjected. Let us now suppose that the system is executing one
of its normal component vibrations, so that w = u, and

u + p2 u = 0 (6),

while w is proportional to another normal function v.

Since x =p/c, we get from (3)

uv dx dy =f e a + dy (7).

The integral on the right is symmetrical with respect to u and v, and
thus

de dy = 0 (8)

where bears the same relation to v that bears to u.

Accordingly, if the normal vibrations represented by u and v have
different periods,

u v dx dy =0 (9).

In obtaining this result, we have made no assumption as to the boundary
conditions beyond what is implied in the absence of reactions against
acceleration, which, if they existed, would appear in the fundamental
equation (3).

If in (8) we suppose ' = , the equation is satisfied identically, and we
cannot infer the value of u2 dx dy. In order to evaluate this integral
we must follow a rather different course.

If u and v be functions satisfying within a certain contour the
equations \(u+ u=0, v+ \kappa' 2 v=0\), we have

(10),

by Green's theorem. Let us now suppose that v is derived from u by
slightly varying , so that
\[ v = u + (du/d \kappa) \kappa, \kappa' = + \kappa,\]
substituting in (10), we find (11);

or, if u vanish on the boundary, (12).

For the application to a circular area of radius r, we have
\(
u=\cos n J,
\)

(r)
\( v=\cos(n) J\),
(r) and thus from (10) on substitution of polar
coordinates and integration with respect to 0,

(x2) F, (er) J, (2) r dr =rJ, (er) J, ('r) - rd, (e r) = ar dy (xr)
(14). Accordingly, if 4 f (dr) 2 T(r) = 2S, (ur) : Ta (er) dr uo dr
ose'' and ' be different, J (eer) T(x) r dr=0 (15),

an equation first proved by Fourier for the case when

Again from (12)

dashes denoting differentiation with respect to r. Now

and thus (16).

This result is general; but if, as in the application to membranes with
fixed boundaries,

then (17).

204. We may use the result just arrived at to simplify the expressions
for T and V. From

we find

whence is derived the normal equation of motion

and a similar equation for . The value of ,, is to be found from the
consideration that etc, denotes the work done by the impressed forces
during a hypothetical displacement c; so that if Z be the impressed
force, reckoned per unit of area,

o= I ZI, Kg) C08 20 1dr dO (5).

These expressions and equations do not apply to the case n= 0, when c
and are amalgamated. We then have (6), (7). As an example, let us
suppose that the initial velocities are zero, and the initial
configuration that assumed under the influence of a constant pressure 7;
thus .

Now by the differential equation, ,

and thus (8);

Substituting this in (7), we see that the initial value of c,, is

For values of n other than zero, and the initial value of c vanish. The
state of the system at time c is expressed by
\[P =  \cos(P)\]
(10),
\[ W = Zag Ty Wel,\]
(11), the summation extending to all the admissible values of

As an example of forced vibrations, we may suppose that Z, still
constant with respect to space, varies as a harmonic function of the
time. This may be taken to represent roughly the circumstances of a
small membrane set in vibration by a train of aerial waves. If Z=cos gt,
we find, nearly as before,

J(r)

The forced vibration is of course independent of 8. It will be seen
that, while none of the symmetrical normal components are missing, their
relative importance may vary greatly, especially if there be a near
approach in value between g and one of the series of quantities p,- df
the approach be very close, the effect of dissipative forces must be
included.

Special problems, 327, 204

[Again, suppose that the force is applied locally at the centre.
By (5) \( \sum \cos(y)\) (13),
if \( Z \cos(qt)\) denote the whole force at time t. From (7)
\[ Z \cos(qt) = pTru. \]
(14) and w is then given by (11). The series is convergent, unless
r = 0.
But this problem would be more naturally attacked by including
in the solutions of (4 § :J0) the second Bessel's function §341.
In this method k=qC and the ratio of constants by which the
two functions of r are multiplied is determined by the boundary
condition. When a coincides with one of the values of
\todo[inline]{missing}
the second function disappears from the solution.]



205. The pitches of the various simple tones and the radii of the nodal
circles depend on the roots of the equation J, (xa) = J, (2) = 0.

If these (exclusive of zero) taken in order of magnitude be called z,'',
2,'' 2,0 Bl cc , then the admissible values of p
are to be found by multiplying the quantities z, by c+a. The particular
solution may then be written

\[w= (2 r) A \cos(n) + B \sin(n) \cos(a)\]

(1), The lowest tone of the group n corresponds to z,"; and since in
this case J, (") does not vanish for any value of r less than a,
there is no interior nodal circle. If we put s =2, J, will vanish, when
r ,'' a = 2," ; 2 that is, when r=a 28 , n
which is the radius of the one interior nodal circle. Similarly if we
take the root z,", we obtain a vibration with s-1 nodal circles
(exclusive of the boundary) whose radii are

All the roots of the equation J, (xa) = =0 are real. For, if possible,
let xa = A+ ip be a root; then Ka=A- tp is also
1 root, and thus by (14) Article 203, .

I J, (xr) J, (tr) r dr = 0.

Now Jd, (r), J, ('r) are conjugate complex quantities, whose product is
necessarily positive ; so that the above equation requires that either
or w vanish. That cannot vanish appears from the consideration that if a
were a pure imaginary, each term of the ascending series for J, would be
positive, and therefore the sum of the series incapable of vanishing. We
conclude that p=0, or that is real'. The same result might be arrived at
from the consideration that only circular functions of the time can
enter into the analytical expression for a normal component vibration.

Dangling
\footnote{ Riemann, p. 260,}


The equation J, (2) =0 has no equal roots (except zero). From equations
(7) and (8) Article 200 we get
I= 2S. Tuy
whence we see that if J,, J,' vanished for the same value of z, J would
also vanish for that value. But in virtue of (8) Article 200 this would
require that all the functions J, vanish for the value
of z in question.
\footnote{Bourget, `\,`Mémoire sur le mouvement vibratoire des membranes
circulaires,'' Ann, de l' école normale, t. XIX., 1866.
In one passage M. Bourget implies that he
has proved that no two Bessels functions of integral order can have the same root,
but 1 cannot find that he has done so. The theorem, however, is probably true
in the case of functions whose orders differ by 1 or 2, it may be easily
proved from the formulas of § 200.}

206. The actual values of z, may be found by interpolation from Hansen's
tables so far as these extend; or formulae may be calculated from the
descending series by the method of successive approximation, expressing
the roots directly. For the important case of the symmetrical vibrations
(n =0), the values of z, may be found from the following, given by Stokes
\footnote{Camb. Phil. Trans. Vol. IX. On the numerical calculation of a class
of definite integrals and infinite series,;
[In accordance with the calculation of Prof. McMahon
the numerator of the last term in (2) has been altered from '245885
to -245270.]}

\todo[inline]{ smudged, exponent of zo }
\[ \frac{z_o^{s}}{\pi} =  s - 0.25 + 0.050661/(4s-1)
-0.053041/(4s-1)^3 +  0.262051/(4s-1)^5  \ldots \]
(1).  For n= 1, the formula is
\[ \frac{z_1^{s}}{\pi} =  s - 0.25 - 0.151982/(4s-1)
+ 0.015399/(4s-1)^3 -  0.245270/(4s-1)^5  \ldots \]
The latter series is convergent enough, even for the first root,
corresponding tos=1. The series (1) will suffice for values of s greater
than unity; but the first root must be calculated independently. The
accompanying table (A) is taken from Stokes' paper, with a slight
difference of notation.

It will be seen either from the formulae, or the table, that the
difference of successive roots of high order is approximately 7. This is
true for all values of n, as is evident from the descending series (10)
Article 200.

page 329, pdf page 343
[The general formula, analogous to (1) and (2), for the roots of
Jn (z) has been investigated by Prof. McMahon. If ni = 4n', and
a = i7r(2M - 1 +4)
(-3).
we have
7 1 - 1 4 (m - 1) (7 - 31)
 8a 3 (Haf
32 0ft - 1)(83™‘ - 982i + 3779)
15(8(1) ‘ ‘
(4).
This formula may be applied not only to integral values of n as in
(1) and (2), but also when n is fractional. The cases of n = and
w = are considered in 207.]

Frequencies of the different simple tones and of the radii of the nodal circles


Table B includes the values of z, which satisfy J, (2), for n =
0, 1,... 5, s=1, 2,... 9.


TABLE A, = for J,(e)=0. Diff for J,(2)=0. Diff

7655 ; ' 1:2197 1:7571 . 2-2330 2-7546 ; 3:2383 37534 . 4-2411 47527 ;
5-2428 57522 ; 6-2439 67519 ; 7:2448 7-7516 ; 8-2454 8-7514 ; 9-24.59
9-7513 . 10-2463

10-7512 ; 11-2466 11-7511 12-2469

When n is considerable the calculation of the earlier roots becomes
troublesome. For very high values of n, z approximates to a ratio of
equality, as may be seen from the consideration that the pitch of the
gravest tone of a very acute sector must tend to coincide with that of a
long parallel strip, whose width is equal to the greatest width of the sector.

206. - NODAL FIGURES. 275

The figures represent the more important normal modes of vibration, and
the numbers affixed give the frequency referred to 18-2
the gravest as unity, together with the radii of the circular nodes
expressed as fractions of the radius of the membrane. In the case of six
nodal diameters the frequency stated is the result of a rough
calculation by myself.

The tones corresponding to the various fundamental modes of the circular
membrane do not belong to a harmonic scale, but there are one or two
approximately harmonic relations which may be worth notice. Thus

x 1594 = 2125 = 2:136 nearly, S x 1594 = 2657 = 2653 nearly, 2 x
1594 = 3:188 = 3156 nearly ;

so that the four gravest modes with nodal diameters only would give a
consonant chord.

The area of the membrane is divided into segments by the nodal system in
such a manner that the sign of the vibration changes whenever a node is
crossed. In those modes of vibration which have nodal diameters there is
evidently no displacement of the centre of inertia of the membrane. In
the case of symmetrical vibrations the displacement of the centre of
inertia is proportional to

i J, (r) r dr = - a (r) +2 Jy (rt r dr = - J, (xa),

an expression which does not vanish for any of the admissible values of
, since J, (z) and J, (z) cannot vanish simultaneously. In all the
symmetrical modes there is therefore a displacement of the centre of
inertia of the membrane.

207. Hitherto we have supposed the circular area of the membrane to be
complete, and the circumference only to be fixed; but it is evident that
our theory virtually includes the solution of other problems, for
example, some cases of a membrane bounded by two concentric circles. The
complete theory for a membrane in the form of a ring requires the second
Bessel's function.

207. FIXED RADIUS. 277

The problem of the membrane in the form of a semi-circle may be regarded
as already solved, since any mode of vibration of which the semi-circle
is capable must be applicable to the
complete circle also. In order to see this, it is only necessary to
attribute to any point in the complementary semi-circle the opposite
motion to that which obtains at its optical image in the bounding
diameter. This line will then require no constraint to keep it nodal.
Similar considerations apply to any sector
whose angle is an aliquot part of two right angles.

When the opening of the sector is arbitrary, the problem may be solved
in terms of Bessel's functions of fractional order. If the fixed radii
are = 0, 0 = 8, the particular solution is

vn

w= PJ. (r) sin (pt-e) 6 (1), a B

where v is an integer. We see that if 8 be an aliquot part of 7, v1 + B
is integral, and the solution is included among those already used for
the complete circle.

An interesting case is when 8 = 27, which corresponds to the problem of
a complete circle, of which the radius 6=0 is constrained to be nodal.

Fig. 38.

We have w=PJ,, (xr) (v0 (pt-e).

When is even, this gives, as might be expected, modes of vibration
possible without the constraint; but, when v is odd, new modes make
their appearance. In fact, in the latter case the descending series for
J terminates, so that the solution is expressible in finite terms. Thus,
when y=1,

w= PD'' gin 36 (pt-c) (2). Jer The values of are given by

(xa=0, or ca=mr.

Thus the circular nodes divide the fixed radius into equal parts, and
the series of tones form a harmonic scale. In the case of the gravest
mode, the whole of the membrane is at any moment deflected on the same
side of its equilibrium position. It is remarkable that the application
of the constraint to the radius 6 =0 makes the problem easier than
before.

If we take y= 3, the solution is

w= PT (A (rr) (86 (pt-e) (3). Fig. 39. In this case the nodal
radii are

and the possible tones are given by the equation CAN HD KO ves ces (4).
To calculate the roots of
\( \tan(a)\) we may assume \(a=(m+)r--y=X--y\), where y
is a positive quantity, which is small when z is large. Substituting
this, we find \( cot(y)= X-y\), whence

=1.1444 This equation is to be solved by successive approximation. It
will readily be found that 13 146

3 = Yi. y=X +5 +7 Xt +g Xt


so that the roots of \(\tan(a)\) are given by

(5), where X=m+4.

In the first quadrant there is no root after zero since \(\tan(x 2\), and in
the second quadrant there is none because the signs of and \(\tan(x)\) are
opposite. The first root after zero is thus in the third quadrant,
corresponding to m=1. Even in this case the series converges:
sufficiently to give the value of the root with considerable accuracy,
while for higher values of m it is all that could be desired. The actual
values of a : mare 14303, 24590, 3:4709, 44747, 54818, 64844, etc.

208. The effect on the periods of a slight inequality in the density of
the circular membrane may be investigated by the general method
Article 90, of which several examples have already been given. It will
be sufficient here to consider the case of a small load M attached to
the membrane at a point whose radius vector is 7'.

We will take first the symmetrical types (n= 0), which may still be
supposed to apply notwithstanding the presence of M. The kinetic energy
T is (6) Article 204 altered from

and therefore (1)

where P denotes the value of p, when there is no load.

The unsymmetrical normal types are not fully determinate for the
unloaded membrane; but for the present purpose they must be taken so as
to make the resulting periods a maximum or minimum, that is to say, so
that the effect of the load is the greatest and least possible. Now,
since a load can never raise the pitch, it is clear that the influence
of the load is the least possible, viz. zero, when the type is such that
a nodal diameter (it is indifferent which) passes through the point at
which the load is attached. The unloaded membrane must be supposed to
have two coincident periods, of which one is unaltered by the addition of the
load. The other type is to be chosen, so that the alteration of period
is as-great as possible, which will evidently be the case when the
radius vector 7' bisects the angle between two adjacent nodal diameters.
Thus, if r' correspond to 6 = 0, we are to take w= Pan In (Kat) 008 08 ;
so that (2) Article 204 vi = t pr a" Pn J,'' (mn) + + Mu P J, (Ka) The
altered p is therefore given by 1 2M TE ant)) p r a J, (Ka) ae e

Of course, if r' be such that the load lies on one of the nodal circles,
neither period is affected.

Pn Pag =

For example, let M be at the centre of the membrane. J, (0) vanishes,
except when n=0; and J,(0)=1. It is only the symmetrical vibrations
whose pitch is influenced by a central load, and for them by (1)

Pad: Pal - (8). By (6) Article 200 SI (e) =-J, (2)

so that the application of the formula requires only a knowledge of the
values of J, (z), when J, (2) vanishes, Article 200. For the gravest
mode the value Of Ty (Hg a) j is `51903'. When x,, is considerable,

FT (High) = 2+ Whey approximately ; so that for the higher components
the influence of the load in altering the pitch increases,

The influence of a small irregularity in disturbing the nodal system may
be calculated from the formulae of Article 90. The most obvious effect
is the breaking up of nodal diameters into curves of hyperbolic form due
to the introduction of subsidiary symmetrical vibrations. In many cases
the disturbance is favoured by close agreement between some of the
natural periods.

209. We will next investigate how the natural vibrations of a uniform
membrane are affected by a slight departure from the exact circular
form.

Dangling
\footnote{The succeeding values are approximately 841, 271, 232, 206, -187, etc.}

Whatever may be the nature of the boundary, w satisfies the equation

(1),

where is a constant to be determined. By Fourier's theorem may be
expanded in the series
\[wo = w + w \cos(0 +4) + 0, \cos(2 A+ a) + + w, \cos(n)+4) + 0,\]
where w,, w,, etc. are functions of r only, Substituting in (1), we see
that w, must satisfy

dw, 1 dw, n ete art (Pp) a0 of which the solution is . J, (xr); for, as
in Article 200, the other function of r cannot appear.

The general expression for w may thus be written y= AT, (er) +d, (er)
(A, 008 0+ B, (0) +... +, (xr) (A (n)+ B, (0) + ... (2). For all points
on the boundary w is to vanish a In the case of a nearly circular
membrane the radius vector is
nearly constant. We may take r=a+6r, dr being a small function of 86.
Hence the boundary condition is

0=A, J, (xa) + cdr J,' (1a) + 2.0008 + J, (a) + 8r J,' (xa) A (nO + B,
(n6 Ti(8),

which is to hold good for all values of 0.

Let us consider first those modes of vibration which are nearly
symmetrical, for which therefore approximately

w= A,J, (kr).

All the remaining coefficients are small relatively to 4,, since the
type of vibration can only differ a little from what it would

a ns

be, were the boundary an exact circle. Hence if the squares of the small
quantities be omitted, (3) becomes

A, J, (wa) + 8r J, (wa) + J, (wa) A (0 + B, (6) +... +d, (xa) A (nO + B,
(n6 +...=0 (4).

If we integrate this equation with respect to 0 between the limits 0 and
27, we obtain

J, (wa) + J,' (1a) =0,

or Jca+ ef" or Fh -9 (5),
which shews that the pitch of the vibration is approximately the same as
if the radius vector had uniformly its mean value.

This result allows us to form a rough estimate of the pitch of any
membrane whose boundary is not extravagantly elongated. If o denote the
area, so that po is the mass of the whole membrane, the frequency of the
gravest tone is approximately

[A numerical error is here corrected.]


In order to investigate the altered type of vibration, we may
multiply (4) by (n), or (n), and then integrate as before. Thus

2

A, J. (xa) i , KBr (nd dd +A, J,(ca) =0

A, J,' (xa) ih cor (nO dO +B, J,,(xa) =0
which determine the ratios A,:.A, and B,: A,.

If r = 6r, + 6r, +... Or, + be Fourier's expansion, the final
expression for w may be written, (8)

When the vibration is not approximately symmetrical, the question
becomes more complicated. The normal modes for the truly circular
membrane are to some extent indeterminate, but the
irregularity in the boundary will, in general, remove the
indeterminateness. The position of the nodal diameters must be taken, so
that the resulting periods may have maximum or minimum values. Let us,
however, suppose that the approximate type is
W= Ay AS, (1) (VO. (9),
and afterwards investigate how the initial line must be taken in order
that this form may hold good.


All the remaining coefficients being treated as small in comparison with
A,, we get from (4)

A,J, (xa) +... + ALA, (ua) + J,' (xa) (v) + B,J, (wa) (vO + + J,
(a) A (nO + B (nO +... =0. (10). - Multiplying by (v) and
integrating, wd, (xa) + ed, a dr (d)d = 0, 0 or

J, xa + w Br (t) =0, 0
which shews that the effective radius of the membrane is

The ratios of A, and B, to A, may be found as-before by integrating
equation (10) after multiplication by (7), (0).

But the point of greatest interest is the pitch. The initial line is to
be so taken as to make the expression (11) a maximum or minimum. If we
refer to a line fixed in space by putting -a instead of 8, we have to
consider the dependence on a of the quantity

oe i dr (v) (6 - a) de, 0 which may also be written

(t) ya i or (v) +2 (va) (vi) i br (v6) (v8)

+ (yz i br sin'' VOD (12),

and is of the form

A (va) + 2B (va) (va) + Cain' va,

A, B, C being independent of a. There are accordingly two admissible
positions for the nodal diameters, one of which makes the period a
maximum, and the other a minimum. The diameters of one set bisect the
angles between the diameters of the other set.

There are, however, cases where the normal modes remain indeterminate,
which happens when the expression (12) is independent of a. This is the
case when 6r is constant, or when 6r is proportional to (v6). For
example, if Sr were proportional to (26), or in other words the boundary
were slightly elliptical, the nodal system corresponding to n=2 (that
consisting of a pair of perpendicular diameters) would be arbitrary in
position, at least to this order of approximation. But the single
diameter, corresponding to n=1, must coincide with one of the principal
axes of the ellipse, and the periods will be different for the two axes.

210. We have seen that the gravest tone of a membrane, whose boundary is
approximately circular, is nearly the same as that of a mechanically
similar membrane in the form of a circle of the same mean radius or
area. If the area of a membrane be given, there must evidently be some
form of boundary for which the pitch (of the principal tone) is the
gravest possible, and this form can be no other than the circle. In the
case of approximate circularity an analytical demonstration may be
given, of which the following is an outline.

The general value of w being
\[ w= A J, (er) +... 4d, (r) (A (n)+ B ) + \]
(1), in which for the present purpose the coefficients \(A\), \(B, \ldots\)
are small relatively to A,, we find from the condition that w vanishes when
r=a+36r,

A, J, (wa) + A, Jy (xa) 8r + 8A,'' (a). (Sr) + , +2 J, (wa) + WJ,' (xa)
dr + 4 (nd) + B (n6) =0... (2). Hence, if br = 2, 008 + 8, (6)+... +a
(n) + 8, (n) + (3),

we obtain on integration with respect to from 0 to 2z r,

from which we see, as before, that if the squares of the small
quantities be neglected, J, (xa) =0, or that to this order of
approximation the mean radius is also the effective radius. In order to
obtain a closer approximation we first determine A, : A, and B,: A, by
multiplying (2) by (n), (n), and then integrating between the limits 0
and 27. Thus

AJ,=-Ka,AJ', BW, =-KB, Ad 0 . (5). Substituting these values in (4), we
get . J,(ea) = 423 a Ge +B) aye -10, (6). Since J, satisfies the
fundamental equation Fe FE TAT, = V ()

and in the present case J,=0 approximately, we may replace

J.'' by -= 2. Equation (6) then becomes

J, (0a) = hit I, mic +00 p+ 1 zat - ...(8). Let us now suppose that a +
da is the equivalent radius of the membrane, so that I, (a+ da) =J, (xa)
+ J, (xa) eda =0. Then by (8) we find

da=- Ge +8 fe +s (9).

Again, if a +da' be the radius of the truly circular membrane of equal
area,

dal = 2-3" (a8 + Bg) (10) 5

= 4 n=l

so that ; iy 2 2 J,' (xa) ae . +B. 1 + 0a eat ose (11). The question is
now as to the sign of the right-hand member. If n=1, and z be written
for xa, J, (2) J, (2) vanishes approximately by (7), since in general
J,=-J,, and in the present case J,(z)=0 nearly. Thus da' -da=0, as
should evidently be the case, since the term in question represents
merely

a displacement of the circle without an alteration in the form of the
boundary. When n= 2, (8) Article 200,

from which and (7) we find that, when J, = 0, (12), whence 2 (13),

which is positive, since z=2,

We have still to prove that

is positive for integral values of n greater than 2, when z= 2-404. For
this purpose we may avail ourselves of a theorem given in Riemann's
Partielle Differential gleichungen, to the effect that neither J, nor
J,' has a root (other than zero) less than n The differential equation
for J, may be put into the form

J = 0;

while initially Jn, and J'n, (as well as d Jn/d log z ) are positive.
Accordingly begins by increasing and does not cease to do so before z =
n, from which it is clear that within the range z= 0 to z=n, neither Jn,
nor \(J_n^{'}\) can vanish. And since J, and J,' are both positive until
z = n, it follows that, when n is an integer greater than 2404, da' - da
is positive. We conclude that, unless a,, ,, a, all vanish, da' is
greater than da, which shews that in the case of any membrane of
approximately circular outline, the circle of equal area exceeds the
circle of equal pitch.

We have seen that a good estimate of the pitch of an approximately
circular membrane may be obtained from its area alone, but by means of
equation (9) a still closer approximation may be effected. We will apply
this method to the case of an ellipse,

-whose semi-axis major is and eccentricity e.

The polar equation of the boundary is r=R l--fei-Get 4+ fe (26) + (14);
so that in the notation of this section a=R(1-je-Ge), a=feR. Accordingly
by (9)

for tan

or by (12), since R =2 = 2404, 2779 Thus the radius of the circle of
equal pitch is

in which the term containing e should be correct.

The result may also be expressed in terms of e and the area o.

We have R= /Z( 1+ i dt ae),

o 3779 , a+da=y/ (1-27) ... (16),

from which we see how small is the influence of a moderate eccentricity,
when the area is given.

and thus

211. When the fixed boundary of a membrane is neither straight nor
circular, the problem of determining its vibrations presents
difficulties which in genera could not be overcome without the
introduction of functions not hitherto discussed or tabulated. A partial
exception must be made in favour of an elliptic boundary ; but for the
purposes of this treatise the importance of the problem is scarcely
sufficient to warrant the introduction of complicated analysis. The
reader is therefore referred to the original investigation of M.
Mathieu
\footnote{Liouville, XIII, 1868; Cours de physique mathematique, 1873, p. 122.}

[The method depends upon the use of conjugate functions. If
\[ x  +  i y =  e \cos ( \xi  + i \eta ) \]
(1), then the level sets of constant \(\eta \) are confocal ellipses,
and of \(\xi\) are confocal hyperbolas.  In terms of \(\xi\), \(\eta\)
the fundamental equation
(V + k ) a = 0 becomes
J + k' (cash - cos f)=0 (2),
where k' = ke.
The solution of (2) may be found in the form
H = (3),
in which S is a function of f only, and H ,a function of 77 only, provided
\[ \frac{d^2}{d \eta^2} \Xi  + (k^{\prime 2} \cosh^2 \xi -a ) \Xi = 0 \]
(4).
\[ \frac{d^2}{d \eta^2} H  + (k^{\prime 2} \cosh^2 \xi -a ) H = 0 \]
(5), \(a\) being an arbitrary constant.
\footnote{Pockels Uber die partielle Differential gliechung \(\Delta u+k^2u=0\), p 114}
Michell
\footnote{Messinger of Mathematics, vol. XIX. p. 86, 1890.}
has shewn that the elliptic transformation (1) is the
only one which yields an equation capable of satisfaction in the form (3).]

Soluble cases may be invented by means of the general solution

w=A,J,(r)++ (A (nd + B, (n6) J, (ar) + 000 For example we might take
w= J, (xr) -dd, (xr) (8)
and attaching different values to A, trace the various forms of boundary
to which the solution will then apply.

Useful information may sometimes be obtained from the theorem of Article 88,
which allows us to prove that any contraction of the fixed boundary
of a vibrating membrane must cause an elevation of pitch, because the
new state of things may be conceived to differ from the old merely by
the introduction of an additional constraint. Springs, without inertia,
are supposed to urge the line of the proposed boundary towards its
equilibrium position, and gradually to become stiffer. At each step the
vibrations become more rapid, until they approach a limit, corresponding
to infinite stiffness of the springs and absolute fixity of their points
of application. It is not necessary that the part cut off should have
the same density as the rest, or even any density at all.


For instance, the pitch of a regular polygon is intermediate between
those of the inscribed and circumscribed circles. Closer limits would
however be obtained by substituting for the circumscribed circle that of
equal area according to the result of Article 210. In the case of the
hexagon, the ratio of the radius of the circle of equal area to that of
the circle inscribed is 1-050, so that the mean
of the two limits cannot differ from the truth by so much as 24 per
cent, In the same way we might conclude that the sector of a circle of
60 is a graver form than the equilateral triangle obtained by
substituting the chord for the arc of the circle. ,


The following table giving the relative frequency in certain calculable
cases for the gravest tone of membranes under similar mechanical
conditions and of equal area (c), shews the effect of a greater or less
departure from the circular form.

\begin{tabular}{lr}
Circle & a = 4.261  \\
Square & S V2.7 = 4.443\\
Quadrant of a circle & 5.135/2 r = 4.551\\
Sector of a circle 60 degrees  & 6.379 Jt = 4.616\\
Rectangle \( 3 \times 2\) &   4.616\\
Equilateral triangle & 2 pi tan 30 degrees = 4.774\\
Semicircle &  3.832  = 4.803\\
Rectangle \( 2 \times 1\) &  \\
Right-angled isosceles  & pi 5/2 = 4.967\\
Rectangle 3 X 1 & pi = 5.736\\
\end{tabular}

For instance, if a square and a circle have the same area, the
former is the more acute in the ratio 4.443 : 4.261

For the circle the absolute frequency is
\[(2 \pi)^{-1} \times \sqrt{\pi}/\sqrt{\sigma},\]
where \( c=\sqrt{T_1}/\sqrt{\rho}\).

In the ease of similar forms the frequency is inversely as the linear
dimension.

[From the principle that an extension of boundary is always
accompanied by a fall of pitch, we may infer that the gravest
mode of a membrane of any shape, and of any variable density, is
devoid of internal nodal lines.]

212. The theory of the free vibrations of a membrane was first
successfully considered by Poisson\footnote{Mém, de l' Académie, t. VIII. 1829}
His theory in the case of the
rectangle left little to be desired, but his treatment
of the circular membrane was restricted to the symmetrical vibrations.
Kirchhoff's solution of the similar, but much more difficult, problem of
the circular plate was published in 1850, and Clebsch's Theory of
Elasticity (1862) gives the general theory of the circular membrane
including the effects of stiffness and of rotatory inertia. It will
therefore be seen that there was not much left to be done in 1866;
nevertheless the memoir of Bourget already referred to contains a useful
discussion of the problem accompanied by very complete numerical
results, the whole of which however were not new.

213.  Observations of M. Bourget.
In his experimental investigations M. Bourget made use of various
materials, of which paper proved to be as good as any, The paper is
immersed in water, and after removal of the superfluous moisture by
blotting paper is placed upon a frame of wood whose edges have been
previously coated with glue. The contraction of the paper in drying
produces the necessary tension, but many failures may be met with before
a satisfactory result is obtained. Even a well stretched membrane
requires considerable precautions in use, being liable to great
variations in pitch in consequence of the varying moisture of the
atmosphere. The vibrations are excited by organ-pipes, of which it is
necessary to have a series proceeding by small intervals of pitch, and
they are made evident to the eye by means of a little sand scattered on
the membrane. If the vibration be sufficiently vigorous, the sand
accumulates on the nodal lines, whose form is thus defined with more or
less precision. Any inequality in the tension shews itself by the
circles becoming elliptic.

The principal results of experiment are the following:

A circular membrane cannot vibrate in unison with every sound. It can
only place itself in unison with sounds more acute than that heard when
the membrane is gently tapped.

As theory indicates, these possible sounds are separated by less'' and
less intervals, the higher they become.

The nodal lines are only formed distinctly in response to certain
definite sounds. A little above or below confusion ensues, and when the
pitch of the pipe is decidedly altered, the membrane remains unmoved.
There is not, as Savart supposed, a continuous transition from one
system of nodal lines to another.

[The reader who wishes to pursue the subject from a mathematical point of
view is referred to an excellent discussion by Pockels (Leipzig, 1891) of the
differential equation 0.]


The nodal lines are circles or diameters or combinations of circles and
diameters, as theory requires. However, when the number of diameters
exceeds two, the sand tends to heap itself confusedly towards the middle
of the membrane, and the nodes are not well defined. ;

The same general laws were verified by MM. Bernard and Bourget in the
case of square membranes'; and these authors consider that the results
of theory are decisively established in opposition to the views of
Savart, who held that a membrane was capable of responding to any sound,
no matter what its pitch might be. But I must here remark that the
distinction between forced and free vibrations does not seem to have
been sufficiently borne in mind. When a membrane is set in motion by
aerial waves having their origin in an organ-pipe, the vibration is
properly speaking forced. Theory asserts, not that the membrane is only
capable of vibrating with certain defined frequencies, but that it is
only capable of so vibrating freely. When however the period of the
force is not approximately equal to one of the natural periods, the
resulting vibration may be insensible.

In Savart's experiments the sound of the pipe was two or three octaves
higher than the gravest tone of the membrane, and was accordingly never
far from unison with one of the series of over tones. MM. Bourget and
Bernard made the experiment under more favourable conditions. When they
sounded a pipe somewhat lower in pitch than the gravest tone of the
membrane, the sand remained at rest, but was thrown into vehement
vibration as unison was approached. So soon as the pipe was decidedly
higher than the membrane, the sand returned again to rest. A
modification of the experiment was made by first tuning a pipe about a
third higher than the membrane when in its natural condition. The
membrane was then heated until its tension had increased sufficiently to
bring the pitch above that of the pipe. During the process of cooling
the pitch gradually fell, and the point of coincidence manifested itself
by the violent motion of the sand, which at the beginning and end of the
experiment was sensibly at rest.

Dangling\footnote{ Ann. de Chimie, UX. 449-479, 1860.}

M. Bourget found a good agreement between theory and observation with
respect to the radii of the circular nodes, though the test was not very
precise, in consequence of the sensible width of
the bands cf sand; but the relative pitch of the various simple tones
deviated considerably from the theoretical estimates. The committee of
the French Academy appointed to report on M. Bourget's memoir suggest as
the explanation the want of perfect fixity of the boundary. It should
also be remembered that the theory proceeds on the supposition of
perfect flexibility-a condition of things not at all closely approached
by an ordinary membrane stretched with a comparatively small force. But
perhaps the most important disturbing cause is the resistance of the
air, which acts with much greater force on a membrane than on a string
or bar in consequence of the large surface exposed. The gravest mode of
vibration, during which the displacement is at all points in the same
direction, might be affected very differently from the higher modes,
which would not require so great a transference of air from one side to
the other.


[In the case of kettle-drums the matter is further complicated by the
action of the shell, which limits the motion of the air upon one side of
the membrane. From the fact that kettle-drums are struck, not in the
centre, but at a point about midway between the centre and edge, we may
infer that the vibrations which it is desired to excite are not of the
symmetrical class. The sound i ; indeed but little affected when the
central point is touched with the finger. Under these circumstances
the principal vibration (1) is that with one nodal diameter and no nodal
circle, and to this corresponds the greater part of the sound obtained
in the normal use of the instrument. Other tones, however, are audible,
which correspond with vibrations characterized (2) by two nodal
diameters and no nodal circle, (3) by three nodal diameters and no nodal
circles, (4) by one nodal diameter and one nodal circle. By observation
with resonators upon a large kettle-drum of 25 inches diameter the pitch
of (2) was found to be about a fifth above (1 ), that of (3) about a
major seventh above (1), and that of (4) a little higher again, forming
an imperfect octave with the principal tone. For the corresponding modes
of a uniform perfectly flexible membrane vibrating in exactly the
theoretical intervals are those represented by the ratios 1'34, 1'66, 1
83 respectively 1-000 : 1-049 : 1175.
In membranes of this kind the tension is due to
capillarity, and is independent of the thickness of the film.
\footnote{The vibrations of soap films have been investigated by Melde. The
frequencies for surfaces of equal area in the form of the Circle, the
square and the equilateral triangle, were found to be J Phil Mag., vol.
IVI., p. 160, 1879. Pogg. Ann., 159, p. 275, 1876. Akustik, p. 131,}
]

\section{Forced vibrations}\label{sect:213a}
213 a.
The forced vibrations of square and circular
membranes have been further experimentally studied by
Elsas \footnote{Nova Acta der Ksl. Leop. Carol. Deutschen Akademie, Bd. XLV. Nr. 1.  Halle, 1882.},
who has confirmed the
conclusions of Savart as to the responsiveness of a membrane to sounds
of arbitrary pitch. In these experiments the vibrations of a fork were
communicated to the membrane by means of a light thread, attached
normally at the centre ; and the position of the nodal curves and of the
maxima of disturbance was traced in the usual manner by sand and
lycopodium. A series of figures accompanies the memoir, shewing the
effect of sounds of progressively rising pitch.

In many instances the curves found do not exhibit the symmetries demanded by
the supposed conditions.  Thus in the case of the square membrane all the
curves should be similarly related to the four corners, and in the case of the
circular membrane all the curves should be circles.
The explanation is probably be sought in the difficulty of
attaining equality of tension. If there be any irregularity, the effect
will be to introduce modes of vibration which should not appear, as
having nodes at the point of excitation, and this especially when there
is a near agreement of periods. Or again, an irregularity may operate to
disturb the balance between two modes of theoretically identical pitch,
which should be excited to the same degree. Indeed the passage through
such a point of isochronism may be expected to be highly unstable in the
absence of moderate dissipative forces.

The theoretical solution of these questions has already (§§ 19C, 204)
been given, but would need much further development for an accurate
determination of the nodal curves relating to periods not included among
the natural periods. But the general course of the phenomenon can be
traced without difficulty.

If the imposed frequency be less than the lowest natural frequency, the
vibration is devoid of (internal) nodes. For a nodal line, if it
existed, being of necessity either endless or terminated at the boundary
\footnote{Otherwise the extremity would have to remain at rest under the action
of component tensions from the surrounding parts which are all in one direction.},
would divide the membrane into two parts. Of these one part would be
vibrating freely with a frequency less than the lowest natural to the
whole membrane, an impossible condition of things (§211). The absence of
nodal curves under the above-mentioned conditions is one of the
conclusions drawn by Elsas from his observations.

As the frequency of the imposed vibration rises through the lowest
natural frequency, a nodal curve manifests itself round the point of
excitation, and gradually extends. The course of things is most easily
followed in the case of the circular membrane, excited at the centre.
The nodal curves are then of necessity also circles, and it is evident
that the first appearance of al nodal circle can take place only at the
centre. Otherwise there would be a circular annulus of finite internal
diameter, vibrating freely with a frequency only infinitesimally higher
than that of the entire circle. At first sight indeed it might appear
that even an infinitely small nodal circle would entail a finite,
elevation of pitch, but a consideration of the solution (§ 204) as
expressed by a combination of Bessel's functions of the first and second
kinds, shews that this is not the case. At the point of isochronism the
second function disappears, and immediately afterwards re-enters with an
infinitely small coefficient. But inasmuch as this function is itself
infinite when r = 0, a nodal circle of vanishing radius is possible.
Accordingly the fixation of the centre of a vibrating circular
membrane does not alter the pitch, a conclusion which may be extended
to the fixation of any number of detached points of a membrane of any
shape.

The effect of gradually increasing frequency upon the nodal system of a
circular membrane may be thus summarized. Below the first proper tone
there is no internal node. As this point is reached, the mode of
vibration identifies itself with the corresponding free mode, and then
an infinitely small nodal circle manifests itself. As the frequency
further increases, this circle expands, until when the second proper
tone is reached, it coincides with the nodal circle of the free
vibration of this frequency. Another infinitely small circle now
appears, and it as well as the first, continually expands, until they
coincide with the nodal system of a free vibration in the third proper
tone. This process continues as the pitch rises, every circle moving
continually outward . At each coincidence with a natural frequency the
nodal system identifies itself with that of the free vibration, and a
new circle begins to form itself at the centre.

The behaviour of a square membrane is of course more difficult to follow
in detail. The transition from Fig. (34) case (4), corresponding to in =
3, = 1, and m = 1, w = 3, to Fig. (36) where vi = 3, = 3, can be
traced in Elsas a curves through such forms as Fig. 39 a.
